\documentclass[10pt,twocolumn]{article}

% ── Geometry ──
\usepackage[a4paper,margin=2cm]{geometry}

% ── Encoding & fonts ──
\usepackage[T1]{fontenc}
\usepackage[utf8]{inputenc}
\usepackage{lmodern}

% ── Maths ──
\usepackage{amsmath,amssymb}

% ── Graphics & tables ──
\usepackage{graphicx}
\usepackage{booktabs}
\usepackage{multirow}
\usepackage{caption}
\usepackage{subcaption}

% ── Colours & hyperlinks ──
\usepackage[dvipsnames]{xcolor}
\usepackage[colorlinks,citecolor=Blue,linkcolor=BrickRed,urlcolor=Plum]{hyperref}

% ── Bibliography ──
\usepackage[numbers,sort&compress]{natbib}

% ── Misc ──
\usepackage{enumitem}
\setlist{nosep,leftmargin=*}

% ── Compact title ──
\usepackage{titling}
\setlength{\droptitle}{-2em}
\pretitle{\begin{center}\Large\bfseries}
\posttitle{\end{center}\vspace{-1em}}
\preauthor{\begin{center}\normalsize}
\postauthor{\end{center}\vspace{-1em}}
\predate{\begin{center}\small}
\postdate{\end{center}\vspace{-1.5em}}

% ── Compact sections ──
\usepackage{titlesec}
\titlespacing*{\section}{0pt}{1.0ex plus .5ex minus .2ex}{0.6ex plus .2ex}
\titlespacing*{\subsection}{0pt}{0.8ex plus .3ex minus .2ex}{0.4ex plus .1ex}

% ══════════════════════════════════════════════════════════════════════════════
\title{Reproducing Fast Pretraining Distillation\\for Small Vision Transformers}
\author{
  Ilyass Moummad \\
  \textit{[University / Program Name]} \\
  \texttt{[email]}
}
\date{\today}

\begin{document}
\maketitle

% ──────────────────────────────────────────────────────────────────────────────
\begin{abstract}
We reproduce and extend the core contributions of TinyViT~\cite{wu2022tinyvit},
a fast distillation framework for small vision transformers.
Using CIFAR-100 and ImageNet-1K, we validate three claims:
(i)~offline logit distillation from a large teacher yields stronger small
models than training from scratch,
(ii)~the benefit transfers across datasets via pretraining then fine-tuning,
and (iii)~teacher architecture choice matters.
As an extension, we propose combining feature-level and logit-level
distillation in an online setting and measure its effect.
All experiments were tracked with Weights~\&~Biases.
\end{abstract}

% ──────────────────────────────────────────────────────────────────────────────
\section{Introduction}
\label{sec:intro}

Vision transformers (ViTs)~\cite{dosovitskiy2020image} achieve state-of-the-art
accuracy on image classification but demand large compute budgets.
Knowledge distillation~\cite{hinton2015distilling} is a natural remedy: a large
\emph{teacher} model transfers its ``dark knowledge'' to a smaller \emph{student}.
However, na\"ive distillation doubles the GPU cost because both models must be
loaded and forwarded simultaneously.

Wu et al.~\cite{wu2022tinyvit} introduce \textbf{TinyViT}, a family of compact
ViTs ($5$--$21$\,M parameters) paired with an \emph{offline distillation}
framework.
The key idea is to \textbf{pre-compute and store} the teacher's top-$K$ logits
for each training image under multiple augmentations, then replay these
sparse logits during student training.
This eliminates the teacher from the training loop entirely, cutting memory
and compute overhead to near zero.

\paragraph{Contributions of the original paper.}
\begin{enumerate}
  \item A new ViT architecture (TinyViT) that combines mobile convolution
        blocks with transformer blocks in a hierarchical design.
  \item An offline distillation framework that stores sparse teacher logits
        once and reuses them across students and epochs.
  \item State-of-the-art accuracy for models under $21$\,M parameters
        ($84.8\%$ top-1 on ImageNet-1K with a $21$\,M model).
\end{enumerate}

\paragraph{Scope of this report.}
We reproduce the distillation framework on CIFAR-100 (direct training) and
ImageNet-1K (pretraining + fine-tuning), compare different teacher
architectures, ablate the logit sparsity parameter~$K$, and extend the
method with feature-level distillation.

% ──────────────────────────────────────────────────────────────────────────────
\section{Background}
\label{sec:background}

\subsection{Offline Logit Distillation}

Given a trained teacher~$T$ and a student~$S$, offline distillation
proceeds in two stages:

\begin{enumerate}
  \item \textbf{Save logits.} For every training image~$x_i$ and $E$ random
        augmentations, compute $z_i^{(e)}=T(x_i^{(e)})$ and store only the
        top-$K$ entries (indices and values) in a compact binary format.
  \item \textbf{Train student.} Replay the same augmentation seeds and train
        $S$ with a soft cross-entropy loss:
        \begin{equation}
          \mathcal{L}_{\text{distill}}
          = -\sum_{c=1}^{C} p_T^{(c)}\;\log\,p_S^{(c)},
          \label{eq:distill}
        \end{equation}
        where $p_T = \mathrm{softmax}(z_T/\tau)$ and
        $p_S = \mathrm{softmax}(z_S/\tau)$ with temperature $\tau$.
\end{enumerate}

Because only the top-$K$ logits are stored (and the rest are imputed as
a uniform residual), the storage footprint for ImageNet-1K at $K\!=\!100$,
$300$~epochs is roughly $16$\,GB.

\subsection{TinyViT Architecture}

TinyViT is a four-stage hybrid architecture:
\begin{itemize}
  \item Stages~1--2 use \textbf{MBConv} (mobile inverted bottleneck) blocks
        for efficient local feature extraction.
  \item Stages~3--4 use \textbf{windowed multi-head self-attention} blocks.
\end{itemize}
Models range from TinyViT-5M ($5.4$\,M params, $1.3$\,GFLOPs) to
TinyViT-21M ($21$\,M params, $4.3$\,GFLOPs).

% ──────────────────────────────────────────────────────────────────────────────
\section{Experimental Setup}
\label{sec:setup}

\subsection{Datasets}

\begin{itemize}
  \item \textbf{CIFAR-100}: $50$k train / $10$k test images, $100$ classes,
        resized to $224\times224$ to match the ViT input resolution.
  \item \textbf{ImageNet-1K}: $1.28$M train / $50$k val images, $1000$
        classes, $224\times224$.
\end{itemize}

\subsection{Models}

Table~\ref{tab:models} lists the models used.
The student is always TinyViT-5M unless stated otherwise.

\begin{table}[h]
\centering
\caption{Models used in our experiments.}
\label{tab:models}
\small
\begin{tabular}{@{}llrl@{}}
\toprule
Model & Type & Params & Role \\
\midrule
TinyViT-5M  & Hybrid ViT & 5.4\,M  & Student \\
TinyViT-21M & Hybrid ViT & 21\,M   & Teacher \\
ViT-Base    & Transformer & 86\,M  & Teacher \\
ResNet-152  & CNN         & 60\,M  & Teacher \\
CLIP-ViT-L/14 & Transformer & 304\,M & Teacher \\
\bottomrule
\end{tabular}
\end{table}

\subsection{Training Details}

All experiments use AdamW with cosine learning rate decay and $10$-epoch
warmup.
CIFAR-100 runs: $200$~epochs, base LR~$10^{-3}$, batch size~$256$.
ImageNet-1K pretraining: $300$~epochs, batch size~$512$ (2~GPUs).
Fine-tuning from ImageNet-1K to CIFAR-100: $30$~epochs, LR~$2.5\times10^{-4}$.
Data augmentation includes RandAugment, random erasing, and label smoothing
($\epsilon=0.1$). Mixup/CutMix are \emph{disabled} during distillation
to match the original paper's protocol.

% ──────────────────────────────────────────────────────────────────────────────
\section{Experiments and Results}
\label{sec:results}

We organise our experiments into four parts.

% ·· Part 1 ··
\subsection{Part 1 --- CIFAR-100 Direct Training}

\paragraph{Goal.} Reproduce TinyViT's offline distillation on CIFAR-100 and
compare multiple teacher architectures.

\paragraph{Protocol.}
(a)~Train TinyViT-5M from scratch (lower bound).
(b)~Fine-tune four teachers (ResNet-152, ViT-Base, TinyViT-21M, CLIP-ViT-L)
on CIFAR-100.
(c)~Save top-$K$ logits for each teacher.
(d)~Train TinyViT-5M students using each teacher's logits.

% ── placeholder figure ──
\begin{figure}[t]
  \centering
  \fbox{\parbox{0.9\linewidth}{\centering
    \vspace{2.5cm}
    \textcolor{gray}{\textit{WandB screenshot: Part~1 validation accuracy curves\\
    (scratch vs.\ distillation from each teacher)}}
    \vspace{0.3cm}
  }}
  \caption{CIFAR-100 validation accuracy for scratch baseline and
           distillation with different teachers (Part~1).}
  \label{fig:part1}
\end{figure}

\begin{table}[t]
\centering
\caption{Part~1 results --- CIFAR-100 direct training with TinyViT-5M.}
\label{tab:part1}
\small
\begin{tabular}{@{}lcc@{}}
\toprule
Experiment & Teacher Acc. & Student Acc. \\
\midrule
Scratch baseline       & ---    & \textit{xx.x\%} \\
\midrule
Distill (TinyViT-21M) & \textit{xx.x\%} & \textit{xx.x\%} \\
Distill (ResNet-152)   & \textit{xx.x\%} & \textit{xx.x\%} \\
Distill (ViT-Base)     & \textit{xx.x\%} & \textit{xx.x\%} \\
Distill (CLIP-ViT-L)  & \textit{xx.x\%} & \textit{xx.x\%} \\
\bottomrule
\end{tabular}
\end{table}

\paragraph{Findings.}
\textit{%
Fill in after completing runs.
Expected: distillation from a stronger teacher (CLIP) should give the
largest improvement; distillation benefit on CIFAR-100 alone is modest
(+2--4\%) because the dataset is small (50k images).}

% ── Top-K ablation ──
\paragraph{Top-$K$ ablation.}
We vary $K \in \{10, 25, 50, 100\}$ when saving logits from both
TinyViT-21M and ViT-Base teachers on CIFAR-100.

\begin{table}[h]
\centering
\caption{Effect of logit sparsity~$K$ on student accuracy (CIFAR-100).}
\label{tab:topk}
\small
\begin{tabular}{@{}ccc@{}}
\toprule
$K$ & TinyViT-21M teacher & ViT-Base teacher \\
\midrule
10  & \textit{xx.x\%} & \textit{xx.x\%} \\
25  & \textit{xx.x\%} & \textit{xx.x\%} \\
50  & \textit{xx.x\%} & \textit{xx.x\%} \\
100 & \textit{xx.x\%} & \textit{xx.x\%} \\
\bottomrule
\end{tabular}
\end{table}

\textit{Expected: $K\!\ge\!25$ should suffice for CIFAR-100 (100 classes),
consistent with the paper's finding that sparse logits retain most of the
information.}

% ·· Part 2 ··
\subsection{Part 2 --- ImageNet-1K Pretraining}

\paragraph{Goal.} Reproduce the paper's core methodology: pretrain
TinyViT on a large-scale dataset (ImageNet-1K) with distillation, then
evaluate by fine-tuning on the target task.

\paragraph{Protocol.}
(a)~Save teacher logits on ImageNet-1K (CLIP-ViT-L via timm, $\sim$88\% top-1).
(b)~Pretrain TinyViT-21M for 300~epochs: \textbf{with} distillation vs.\
\textbf{without} (scratch baseline).

\begin{figure}[t]
  \centering
  \fbox{\parbox{0.9\linewidth}{\centering
    \vspace{2.5cm}
    \textcolor{gray}{\textit{WandB screenshot: ImageNet-1K pretraining curves\\
    (distillation vs.\ scratch for TinyViT-21M)}}
    \vspace{0.3cm}
  }}
  \caption{ImageNet-1K validation accuracy during pretraining (Part~2).}
  \label{fig:part2}
\end{figure}

\begin{table}[h]
\centering
\caption{Part~2 results --- ImageNet-1K pretraining.}
\label{tab:part2}
\small
\begin{tabular}{@{}lcc@{}}
\toprule
Pretraining & IN-1K Val Acc. & Paper Ref. \\
\midrule
TinyViT-21M scratch      & \textit{xx.x\%} & 83.1\% \\
TinyViT-21M + distill (CLIP) & \textit{xx.x\%} & 84.8\%$^*$ \\
\bottomrule
\multicolumn{3}{@{}l}{\footnotesize $^*$Paper result is with IN-22K distillation.}
\end{tabular}
\end{table}

\paragraph{Findings.}
\textit{Fill in.
Note: we use IN-1K (1.28\,M images) instead of IN-22K (14\,M),
so a gap relative to the paper's numbers is expected.}

% ·· Part 3 ··
\subsection{Part 3 --- Transfer Learning (IN-1K $\to$ CIFAR-100)}

\paragraph{Goal.} Measure how much distillation \emph{during pretraining}
benefits downstream fine-tuning on a different dataset.

\paragraph{Protocol.}
Fine-tune the two pretrained TinyViT-21M checkpoints from Part~2 on CIFAR-100
for $30$ epochs with identical hyperparameters.

\begin{table}[h]
\centering
\caption{Part~3 results --- Transfer to CIFAR-100.}
\label{tab:part3}
\small
\begin{tabular}{@{}lcc@{}}
\toprule
Pretraining & CIFAR-100 Acc. & $\Delta$ \\
\midrule
Scratch $\to$ fine-tune  & \textit{xx.x\%} & --- \\
Distill $\to$ fine-tune  & \textit{xx.x\%} & \textit{+x.x} \\
\bottomrule
\end{tabular}
\end{table}

\begin{figure}[h]
  \centering
  \fbox{\parbox{0.9\linewidth}{\centering
    \vspace{2.5cm}
    \textcolor{gray}{\textit{WandB screenshot: CIFAR-100 fine-tuning curves\\
    (from scratch vs.\ from distilled pretrained)}}
    \vspace{0.3cm}
  }}
  \caption{CIFAR-100 fine-tuning from scratch vs.\ distilled pretraining (Part~3).}
  \label{fig:part3}
\end{figure}

\paragraph{Findings.}
\textit{Fill in.
Expected: the distillation-pretrained model should converge faster and
achieve 2--3\% higher final accuracy, confirming that distillation
benefits transfer across datasets.}

% ·· Part 4 ··
\subsection{Part 4 --- Feature Distillation Extension}

\paragraph{Goal.}
Extend the TinyViT framework by adding \textbf{feature-level} distillation
alongside logit-level distillation, in a fully \emph{online} setting
(no saved logits required).

\paragraph{Method.}
The total loss is:
\begin{equation}
  \mathcal{L} = \mathcal{L}_{\text{KL}} + \beta\;\mathcal{L}_{\text{feat}},
  \label{eq:feature}
\end{equation}
where $\mathcal{L}_{\text{KL}}$ is the KL divergence between student and
teacher logits (computed online), and
\begin{equation}
  \mathcal{L}_{\text{feat}} = 1 - \cos\!\bigl(
    \mathrm{proj}(\mathbf{f}_S),\;\mathbf{f}_T
  \bigr)
\end{equation}
with a learnable linear projection from the student feature dimension
($320$ for TinyViT-5M) to the teacher dimension ($576$ for TinyViT-21M).
We set $\beta=0.5$ by default.

\begin{table}[h]
\centering
\caption{Part~4 results --- Online distillation on CIFAR-100.}
\label{tab:part4}
\small
\begin{tabular}{@{}lc@{}}
\toprule
Method & CIFAR-100 Acc. \\
\midrule
Scratch (no distillation) & \textit{xx.x\%} \\
Online logits only        & \textit{xx.x\%} \\
Online logits + features ($\beta\!=\!0.5$) & \textit{xx.x\%} \\
\bottomrule
\end{tabular}
\end{table}

\begin{figure}[h]
  \centering
  \fbox{\parbox{0.9\linewidth}{\centering
    \vspace{2.5cm}
    \textcolor{gray}{\textit{WandB screenshot: Online distillation\\
    (logits-only vs.\ logits+features)}}
    \vspace{0.3cm}
  }}
  \caption{Feature distillation extension --- training curves (Part~4).}
  \label{fig:part4}
\end{figure}

\paragraph{Findings.}
\textit{Fill in.
Expected: adding feature distillation should provide a modest improvement
over logits-only online distillation, by also aligning intermediate
representations.}

% ──────────────────────────────────────────────────────────────────────────────
\section{Discussion}
\label{sec:discussion}

\paragraph{Offline vs.\ online distillation.}
The offline approach (Parts~1--3) trades storage for compute:
once teacher logits are saved, student training runs at the cost of
normal training.
The online approach (Part~4) avoids storage but doubles the forward-pass
cost.
On a small dataset like CIFAR-100 ($50$k images), the storage overhead is
negligible ($<100$\,MB), making offline distillation clearly advantageous.
On ImageNet-1K, the storage cost ($\sim$16\,GB for $K\!=\!100$, $300$
epochs) remains practical.

\paragraph{Teacher architecture matters.}
\textit{Discuss the relative rankings of ResNet-152, ViT-Base, TinyViT-21M,
and CLIP-ViT-L as teachers.  Stronger teachers (CLIP) should lead to better
students, but there may be diminishing returns.}

\paragraph{Top-$K$ sparsity.}
\textit{Discuss the top-$K$ ablation.  On CIFAR-100 ($C\!=\!100$),
$K\!=\!25$ already captures most information since $K/C = 25\%$.  On
ImageNet-1K ($C\!=\!1000$), $K\!=\!100$ ($10\%$) is the paper's default.}

\paragraph{Limitations.}
Our reproduction uses ImageNet-1K (1.28\,M images) rather than
ImageNet-22K (14\,M), so we cannot directly match the paper's headline
results.
CIFAR-100 is also a relatively small benchmark; the distillation
advantage is expected to be more pronounced on larger-scale datasets.

% ──────────────────────────────────────────────────────────────────────────────
\section{Conclusion}
\label{sec:conclusion}

We successfully reproduced TinyViT's offline distillation pipeline on
CIFAR-100 and ImageNet-1K, confirming that:
\begin{enumerate}
  \item Sparse teacher logits can replace the full teacher during training
        with negligible loss of information.
  \item Distillation pretraining transfers to downstream tasks, improving
        CIFAR-100 accuracy over a scratch-trained baseline.
  \item Teacher quality positively correlates with student performance.
\end{enumerate}
As an extension, we showed that combining feature and logit distillation
in an online setting provides additional benefits, at the cost of
higher training-time compute.

% ──────────────────────────────────────────────────────────────────────────────
\bibliographystyle{unsrtnat}
\begin{thebibliography}{9}

\bibitem{wu2022tinyvit}
K.~Wu, J.~Zhang, H.~Peng, M.~Liu, B.~Xiao, J.~Fu, and L.~Yuan,
``TinyViT: Fast Pretraining Distillation for Small Vision Transformers,''
in \emph{European Conference on Computer Vision (ECCV)}, 2022.

\bibitem{dosovitskiy2020image}
A.~Dosovitskiy, L.~Beyer, A.~Kolesnikov, D.~Weissenborn, X.~Zhai,
T.~Unterthiner, M.~Dehghani, M.~Minderer, G.~Heigold, S.~Gelly, J.~Uszkoreit,
and N.~Houlsby,
``An Image is Worth 16x16 Words: Transformers for Image Recognition at Scale,''
in \emph{International Conference on Learning Representations (ICLR)}, 2021.

\bibitem{hinton2015distilling}
G.~Hinton, O.~Vinyals, and J.~Dean,
``Distilling the Knowledge in a Neural Network,''
\emph{arXiv:1503.02531}, 2015.

\bibitem{radford2021clip}
A.~Radford, J.~W.~Kim, C.~Hallacy, A.~Ramesh, G.~Goh, S.~Agarwal,
G.~Sastry, A.~Askell, P.~Mishkin, J.~Clark, G.~Krueger, and I.~Sutskever,
``Learning Transferable Visual Models From Natural Language Supervision,''
in \emph{International Conference on Machine Learning (ICML)}, 2021.

\bibitem{touvron2021deit}
H.~Touvron, M.~Cord, M.~Douze, F.~Massa, A.~Sablayrolles, and H.~J{\'e}gou,
``Training data-efficient image transformers \& distillation through attention,''
in \emph{International Conference on Machine Learning (ICML)}, 2021.

\end{thebibliography}

\end{document}
